\documentclass{amsart}

% 数学相关宏包
\usepackage{amsmath,mathtools,amsthm,amssymb}
\usepackage{bbm}
\usepackage{mathrsfs}
\usepackage{centernot}

% 图形和绘图
\usepackage{graphicx,float}
\usepackage{tikz,tikz-cd}
\usepackage{pgfplots}
\usepackage[framemethod=TikZ]{mdframed}
\usetikzlibrary{decorations.pathmorphing,positioning,arrows,automata,calc,shapes,patterns}
\pgfplotsset{compat=1.18}
\usepgfplotslibrary{statistics}

% 表格和列表
\usepackage{booktabs,multirow,colortbl}
\usepackage{enumitem}
\setlist[itemize,enumerate,description]{noitemsep}
\setlistdepth{9}

% 算法
\usepackage{algorithm,algorithmic}

% 字体和颜色
\usepackage{xcolor,listings}
\usepackage{xeCJK}
\setCJKmainfont{Noto Serif CJK SC}

% 页面和链接
\usepackage{fancyhdr,caption}
\usepackage{hyperref}
\usepackage{cleveref}

% 工具宏包
\usepackage{etoolbox,letltxmacro,docmute}
\usepackage{gensymb}

% 自定义设置
\renewcommand{\arraystretch}{1.5}
\let\originalmathbb\mathbb
\renewcommand{\mathbb}[1]{%
    \ifstrequal{#1}{1}{\mathbbm{#1}}{\originalmathbb{#1}}%
}

% 颜色定义
\definecolor{lightblue}{RGB}{173,216,230}
\definecolor{lightgreen}{RGB}{144,238,144}
\definecolor{lightyellow}{RGB}{255,255,224}
\definecolor{lightgray}{RGB}{211,211,211}

		% 超链接设置
		\hypersetup{
		    colorlinks=true,
		    linkcolor=red,
		    filecolor=magenta,
		    urlcolor=cyan,
		    pdftitle={\@title},
		    pdfauthor={\@author},
		    pdfsubject={数学笔记},
		    pdfkeywords={数学},
		    bookmarksnumbered=true,
		    bookmarksopen=true,
		    bookmarksopenlevel=6,
		    pdfstartview=FitH,
		    pdfpagemode=UseOutlines,
		    pdfborder={0 0 0}
		}

% 定理环境
\theoremstyle{plain}
\newtheorem{theorem}{Theorem}[section]
\newtheorem{lemma}[theorem]{Lemma}
\newtheorem{corollary}[theorem]{Corollary}
\newtheorem{proposition}[theorem]{Proposition}

\theoremstyle{definition}
\newtheorem{definition}[theorem]{Definition}
\newtheorem{example}[theorem]{Example}
\newtheorem{exercise}[theorem]{Exercise}
\newtheorem{problem}[theorem]{Problem}

\theoremstyle{remark}
\newtheorem{remark}[theorem]{Remark}
\newtheorem{note}[theorem]{Note}
\newtheorem{observation}[theorem]{Observation}

% 解答环境
\newtheorem*{solution}{Solution}
\newtheorem*{proof*}{Proof}

% 图片路径
\graphicspath{{F:/资料/2025-summer/figures/}}

\author{乐绎华, Yue Yihua}
\address{中山大学, Sun Yat-sen University}
\email{yueyh@mail2.sysu.edu.cn}
\date{\today}
\title{Mathematical Notes}

\begin{document}

\maketitle
\tableofcontents

For the simplex, the exponential representation is
\[
(\Delta _n,\oplus ;\mathbb{R},\odot )\to \mathrm{Hom}(\Omega^{e}_{x};\mathbb{R})
\]
\[
p\mapsto p^{(e)}(\cdot \mid x)
\]
\[
p\oplus q\mapsto p^{(e)}(\cdot \mid x+y)
\]
\[
c\odot p\mapsto p^{(e)}(\cdot \mid cx)
\]
where $x=\left( \log\frac{p_1}{p_0},\dots,\log\frac{p_n}{p_0} \right)$.

Consider the dual space of $\Delta _n$, denoted $\Delta _n^{*}$, defined in the usual way as
\[
\Delta _n^{*}:=\mathrm{Hom}(\Delta _n,\mathbb{R}),
\]
the set of all real‐valued linear functionals with respect to the Aitchison vector‐space structure on $\Delta_n$.  For $f,g\in\Delta _n^{*}$ and $c\in\mathbb{R}$ we endow $\Delta _n^{*}$ with the \textit{ordinary} pointwise operations
\[
(f+g)(p)\coloneqq f(p)+g(p),\qquad (c\,f)(p)\coloneqq c\,f(p).
\]
The relation between the exponential representation $p^{(e)}(\cdot \mid x)$ and mixture representation $p^{(m)}(\cdot \mid u)$ is
\[
u_i=\frac{e^{ x^{i} }}{1+\sum_{i=1}^{n} e^{ x^{i} }}\iff x^{i}=\log\frac{u_i}{1-\sum_{i=1}^{n} u_i}
\]
Then, for every $p\in\Delta_n$,
\[
(f+g)(p)=f(p)+g(p),\qquad (c\,f)(p)=c\,f(p).
\]
In particular, if we encode the value of a functional through its \textbf{mixture parameters}
$f(p)=p^{(m)}(\cdot\mid u^{(f)}_p)$ and $g(p)=p^{(m)}(\cdot\mid u^{(g)}_p)$, then
\[
(f+g)(p)\mapsto p^{(m)}\bigl(\cdot\mid u^{(f)}_p+u^{(g)}_p\bigr),\qquad
(c\,f)(p)\mapsto p^{(m)}\bigl(\cdot\mid c\,u^{(f)}_p\bigr).
\]
This shows that the \textbf{ordinary} vector‐space structure on the dual space is compatible with the mixture parameterisation.

The centered log‐ratio (clr) of $p$ is itself a functional and hence belongs to $\Delta _n^{*}$.

An inner product in $\Delta _n$ is defined by
\[
\left< p,q \right>  =\sum_{i=0}^{n} \text{clr}_i(p)\text{clr}_i(q)=\frac{1}{2(n+1)}\sum_{i=0}^{n} \sum_{j=0}^{n} \log\frac{p_i}{p_j}\log\frac{q_i}{q_j}
\]
It induces a norm in $\Delta _n$,
\[
\lVert p \rVert =\frac{1}{2(n+1)}\sum_{i=0}^{n} \sum_{j=0}^{n} \left( \log\frac{p_i}{p_j} \right)^{2}\qquad \text{for }p\in \Delta _n
\]
The norm above is an example denoted by $L^{2}$ -norm. In general, for any norm in $\Delta _n$, its dual norm in $\Delta _n^{*}$ is well-defined by
\[
\lVert f \rVert _{*}=\sup_{\lVert p \rVert =1}\lVert f(p) \rVert =\sup_{p\in\Delta _n}\frac{\lVert f(p) \rVert }{\lVert p \rVert }\qquad \text{for }f\in \Delta _n^{*}
\]
It must satisfy the standard norm axioms

\begin{itemize}
	\item $\lVert f+g \rVert_{*}\leq\lVert f \rVert_{*}+\lVert g \rVert_{*}$ (triangle inequality),
	\item $\lVert c\,f \rVert_{*}=|c|\,\lVert f \rVert_{*}$ (positive homogeneity),
	\item $\lVert f \rVert_{*}=0\iff f=\mathbf{o}$.
\end{itemize}

The exponential parameterization $x=[x^{1},\dots,x^{n}]$ has the advantage to express the vector space structure of $\Delta _n$ readily. For $\mathrm{p}=p(\zeta \mid x), \mathrm{q}=p(\zeta \mid y)$, it can be shown that

\begin{enumerate}
	\item $\mathrm{p} \oplus \mathrm{q}=p(\zeta \mid x+y)$;
	\item $c \odot \mathrm{p}=p(\zeta \mid c x)$.
\end{enumerate}

We wonder if the mixture parameterization $u=[u_1,\dots,u_n]$ also has advantage to express the vector space structure of $\Delta _n^{*}$ for some specific norm over $\Delta _n^{*}$, since the mixture parameterization is the dual of exponential parameteriazation.


\end{document}